\documentclass[12pt,a4paper]{article}

\usepackage[UTF8]{ctex}
\usepackage{geometry}
\usepackage{amsmath}
\usepackage{amssymb}
\usepackage{listings}
\usepackage{xcolor}
\usepackage{hyperref}
\usepackage{booktabs}
\usepackage{tabularx}
\usepackage{enumitem}
\usepackage{float}
\usepackage{fancyhdr}
\usepackage{titlesec}

\geometry{left=2.5cm,right=2.5cm,top=3cm,bottom=3cm}
\setlength{\headheight}{15pt}

\lstset{
    basicstyle=\ttfamily\small,
    keywordstyle=\color{blue}\bfseries,
    commentstyle=\color{gray},
    stringstyle=\color{red},
    numbers=left,
    numberstyle=\tiny\color{gray},
    frame=single,
    breaklines=true,
    tabsize=4,
    showstringspaces=false
}

\hypersetup{
    colorlinks=true,
    linkcolor=blue,
    filecolor=magenta,
    urlcolor=cyan,
    pdftitle={AUV ZIT6 Firmware System Documentation}
}

\pagestyle{fancy}
\fancyhf{}
\fancyhead[L]{AUV ZIT6 固件系统}
\fancyhead[R]{\thepage}
\renewcommand{\headrulewidth}{0.4pt}

\titleformat{\section}{\Large\bfseries}{\thesection}{1em}{}
\titleformat{\subsection}{\large\bfseries}{\thesubsection}{1em}{}
\titleformat{\subsubsection}{\normalsize\bfseries}{\thesubsubsection}{1em}{}

\title{\textbf{AUV ZIT6 固件系统}\\ \Large 详细设计与实现文档}
\author{ZIT6 开发团队}
\date{\today}

\begin{document}

\maketitle
\newpage

\tableofcontents
\newpage

% ============================================================================
\section{项目概述}

\subsection{项目简介}

AUV ZIT6 是一款基于 STM32H7 系列微控制器的自主水下航行器（Autonomous Underwater Vehicle, AUV）固件系统。项目采用现代 C++17 模块化架构设计，实现了 6-DOF（六自由度）内部状态/级联控制框架，支持位置环、速度环和直接推力控制三种工作模式。上位机 setpoint 保留 6-DOF 载荷用于统一接口和姿态观测，但 Roll/Pitch 外部设定值始终旁路，现有推力矩阵不变。

\subsection{硬件平台}

\begin{itemize}
    \item \textbf{微控制器}：STM32H743ZI（ARM Cortex-M7，双精度 FPU）
    \item \textbf{通信外设}：
    \begin{itemize}
        \item USART2 — micro-ROS（DDS 协议，与上位机通信）
        \item UART6 — VIT6 下位机（推进器控制板）
        \item UART7 — NAV-300 惯导
        \item USART1 — USBL 超短基线定位系统
        \item I2C1/2 — MS5837 深度传感器
    \end{itemize}
    \item \textbf{传感器}：NAV-300 惯导（INS）、MS5837 压力传感器、DVL 多普勒测速仪
    \item \textbf{操作系统}：FreeRTOS（CMSIS-RTOS V2 封装）
\end{itemize}

\subsection{软件栈}

\begin{table}[H]
\centering
\begin{tabularx}{\textwidth}{lX}
\toprule
\textbf{组件} & \textbf{技术选型} \\
\midrule
微控制器 & STM32H743ZI（ARM Cortex-M7） \\
RTOS & FreeRTOS + CMSIS-RTOS V2 \\
通信 & micro-ROS（基于 DDS/UXR-CE） \\
编程语言 & C++17，部分 C \\
数值计算 & Eigen 3.4+（仅 Core，零动态内存） \\
构建系统 & CMake 3.22+（STM32CubeMX 集成） \\
调试 & SEGGER RTT / J-Link / UART \\
推进器协议 & 定长帧 FA AF [ID] [Payload] FB BF \\
\bottomrule
\end{tabularx}
\caption{软件栈概览}
\end{table}

\subsection{目录结构}

\begin{lstlisting}[caption={项目顶层目录结构}]
UserApp/
├── Algorithm/           # 纯数值算法（无 RTOS 依赖）
├── Application/         # FreeRTOS 任务层
├── Common/             # 全局上下文与工具类
├── Component/          # 高内聚功能组件
│   ├── Chassis/        # 底盘控制
│   ├── ConfigService/  # 参数配置服务
│   ├── HitlSimulator/  # HIL 仿真器
│   ├── MicroRos/       # micro-ROS 封装
│   ├── RosLogger/      # 日志队列
│   └── Safety/         # 安全监控
├── Peripherals/        # 外设驱动与 HAL
├── Porting/            # 硬件适配层
└── Config/             # 配置管理
\end{lstlisting}

% ============================================================================
\section{Safety — 安全监控}

\subsection{SafetyMonitor}

独立安全监控模块，从 ControlTask 中抽离，在 MonitorTask 中以 10Hz 运行。

\subsubsection{ARM / DISARM 逻辑}

\begin{table}[H]
\centering
\begin{tabularx}{\textwidth}{lX}
\toprule
\textbf{参数} & \textbf{值} \\
\midrule
最小解锁心跳数 & 10 次 \\
最小解锁持续时间 & 1000 ms \\
建议心跳频率 & 2 Hz \\
已解锁心跳超时 & 1000 ms \\
未解锁心跳超时清零 & 1000 ms \\
PWM 输出使能 & 推力 $[-1.0, 1.0]$ 映射为 $[-1000, 1000]$ \\
\bottomrule
\end{tabularx}
\caption{安全监控参数}
\end{table}

\subsubsection{解锁条件}
\begin{enumerate}
    \item 心跳计数 $\geq$ 10 且持续 $\geq$ 1000 ms
    \item \texttt{last\_heartbeat\_data == 1} + 导航有效，或 \texttt{data == 3} 跳过导航检查
    \item 解锁时自动记录 Home Offset（当前位置为原点）
\end{enumerate}

\subsubsection{上锁条件}
\begin{enumerate}
    \item 已解锁状态下 500 ms 内未收到心跳 $\rightarrow$ 强制上锁
    \item 上锁时清空 Home Offset，控制层级切回 NONE
\end{enumerate}

\subsection{SoftWatchdog}

软件看门狗，监控三个组件的心跳状态：

\begin{itemize}
    \item \texttt{MICROROS} — micro-ROS 通信状态
    \item \texttt{INS} — 惯导数据接收状态
    \item \texttt{DEPTH} — 深度传感器数据状态
\end{itemize}

所有组件正常时，MonitorTask 喂硬件看门狗（IWDG）。

% ============================================================================
\section{Peripherals — 外设驱动}

\subsection{INS\_Driver（NAV-300 惯导）}

采用函数指针接口解耦硬件操作：

\begin{itemize}
    \item \texttt{InsPortOps} 结构体封装 \texttt{init/read/transmit} 三个函数指针
    \item 字节流解析，帧长 133-256 字节
    \item 输出 \texttt{MotionContext::NavState}（6-DOF 位姿 + 6-DOF 速度）
    \item 支持命令：DVL 开关、重启、位置重置、经纬度装订
\end{itemize}

\subsection{MotionController\_Driver（动力控制板）}

通过 VIT6 下位机通信协议控制执行器：

\begin{lstlisting}[caption={协议数据包格式}]
struct ThrustPacket {
    uint8_t head[2];   // FA AF
    uint8_t id;        // 0x01(推力)
    float Fx, Fy, Fz;  // 三轴力 (N)
    float Fyaw, Fpitch, Froll;  // 三轴力矩
    uint8_t tail[2];   // FB BF
};
\end{lstlisting}

支持命令 ID：\texttt{0x01（推力）}、\texttt{0x02（舵机）}、\texttt{0x03（灯光）}、\texttt{0x07（推力曲线）}

\subsection{USBL\_Driver（超短基线定位）}

通过 USART1 + DMA 接收 USBL 数据：
\begin{itemize}
    \item 帧头 \texttt{EB BE}，帧尾 \texttt{ED DE}
    \item 133 字节定长帧，异或校验
    \item 输出：信标北向/东向/深度、基阵姿态、信号能量
\end{itemize}

\subsection{MS5837\_Driver（深度传感器）}

支持 I2C 和 UART 两种通信后端（通过 DepthCalcBoard 解算板），提供深度和温度测量。

\subsection{External Drivers}
\begin{itemize}
    \item \textbf{HAL}：STM32CubeMX 生成的 HAL 驱动（I2C、UART、DMA、IWDG、TIM）
    \item \textbf{FreeRTOS}：CMSIS-RTOS V2 封装，heap\_4 内存管理
\end{itemize}

% ============================================================================
\section{Porting — 硬件适配层}

Porting 层封装具体硬件引脚和中断服务，提供 \textbf{函数指针接口表} 供上层驱动使用。

\subsection{函数指针解耦模式}

驱动类（如 INS\_Driver）的构造函数接收 \texttt{PortOps} 结构体，运行时绑定具体硬件实现：

\begin{lstlisting}[caption={PortOps 接口模式示例}]
struct InsPortOps {
    void *ctx;
    bool (*init)(void *ctx);
    uint16_t (*read)(void *ctx, uint8_t *buf, uint16_t max_len);
    bool (*transmit)(void *ctx, const uint8_t *data, uint16_t len);
};

struct MotorPortOps {
    void *ctx;
    bool (*transmitDMA)(void *ctx, const uint8_t *data, uint16_t size);
    void *(*getTxPacket)(void *ctx);
};
\end{lstlisting}

\subsection{Porting 实现}

\begin{itemize}
    \item \textbf{INS\_Porting}：UART7 DMA 实现 INS\_Driver 的 InsPortOps
    \item \textbf{MotionController\_Porting}：UART6 DMA 实现 MotionController\_Driver 的 MotorPortOps
    \item \textbf{MS5837\_Porting}：I2C/UART 实现深度传感器后端
    \item \textbf{DepthCalcBoard\_Porting}：深度解算板通信适配
\end{itemize}

% ============================================================================
\section{Common — 全局上下文}

\subsection{AppContext — 依赖注入容器}

使用结构体聚合所有驱动和组件指针，通过构造函数注入任务：

\begin{lstlisting}[caption={AppContext 定义}]
struct AppContext {
    INS_Driver *ins_driver;
    MotionController_Driver *motor_driver;
    MS5837_Driver *depth_sensor;
    RosLogger *logger;
    SoftWatchdog *watchdog;
    ChassisManager *chassis;
};
extern AppContext g_app_ctx;
\end{lstlisting}

\subsection{MotionContext — 运动状态共享}

跨任务共享的运动状态，所有字段使用 \texttt{LockedField<T>} 临界区保护：
\begin{itemize}
    \item \texttt{nav\_state\_} — 导航状态（6-DOF 位姿 + 6-DOF 速度）
    \item \texttt{current\_setpoint\_} — 当前目标设定值
    \item \texttt{last\_output\_forces\_} — 6-DOF 输出力快照
    \item \texttt{home\_offset\_} — 解锁原点偏移
\end{itemize}

\subsection{SystemContext — 系统状态}

\begin{itemize}
    \item \texttt{arm\_state\_} — 解锁/上锁状态 + 心跳计数
    \item \texttt{nav\_status\_} — 惯导/DVL 状态标志
    \item \texttt{is\_planner\_active} — 规划器启用标志
\end{itemize}

\subsection{LockedField — 线程安全字段}

模板类，对任意可拷贝类型 T 自动加 \texttt{taskENTER\_CRITICAL()} / \texttt{taskEXIT\_CRITICAL()}：
\begin{itemize}
    \item \texttt{get()} — 临界区读取
    \item \texttt{set(v)} — 临界区写入
\end{itemize}

% ============================================================================
\section{Application — 任务层}

\subsection{ControlTask（100Hz 控制主循环）}

\begin{lstlisting}[caption={ControlTask 主循环}]
void ControlTask::run() {
    init();
    for (;;) {
        updateNavigation();    // 1. 获取 INS / SITL / HIL 导航数据
        computeAndPublish();    // 2. 执行底盘控制 + 推力输出
        vTaskDelayUntil(&last_wake_time_, pdMS_TO_TICKS(10));
    }
}
\end{lstlisting}

\texttt{updateNavigation()} 数据源选择：
\begin{itemize}
    \item HITL 模式：由 \texttt{HitlSimulator} 提供
    \item SITL 模式：由 \texttt{/zit6/sim/nav} 话题队列提供
    \item 实物模式：\texttt{INS\_Driver::getNavState()}，并融合 MS5837 深度
\end{itemize}

\subsection{MicroRosTask（通信任务）}

封装 micro-ROS 初始化、发布、订阅和服务：

\begin{lstlisting}[caption={MicroRosTask 结构}]
class MicroRosTask {
    MicroRosTransport transport_;
    MicroRosPublisher  publisher_;
    MicroRosSubscriber subscriber_;
    MicroRosService    service_;
};
\end{lstlisting}

\subsection{MonitorTask（10Hz 安全监控）}

\begin{lstlisting}[caption={MonitorTask 主循环}]
void MonitorTask::run() {
    init();
    for (;;) {
        refreshHardwareWatchdogIfNeeded();  // 软狗检查 $\rightarrow$ 喂硬狗
        safety_monitor_.check(HAL_GetTick()); // ARM/DISARM 状态机
        vTaskDelayUntil(&last_wake_time_, pdMS_TO_TICKS(100));
    }
}
\end{lstlisting}

% ============================================================================
\section{通信协议}

\subsection{控制指令 — ZitSetpoint}

\subsubsection{消息结构}

\begin{lstlisting}[caption={ZitSetpoint 消息定义}]
# zit6_interfaces/msg/ZitSetpoint.msg
uint8 control_key   # 控制模式 + 标志位
uint8 type_mask     # 轴掩码
float32 x y z       # 目标值（m/m/s/N）
float32 roll        # 横滚（rad/rad/s）
float32 pitch       # 俯仰（rad/rad/s）
float32 yaw         # 偏航（rad/rad/s）
uint32 seq          # 序列号（可选）
\end{lstlisting}

\subsubsection{control\_key}

\begin{table}[H]
\centering
\begin{tabularx}{\textwidth}{lX}
\toprule
\textbf{位域} & \textbf{含义} \\
\midrule
Bit 0-1 & 模式：0=POSITION, 1=VELOCITY, 2=FORCE \\
Bit 4 (0x10) & 机体系坐标（默认世界系 NED） \\
Bit 5 (0x20) & 增量模式 \\
\bottomrule
\end{tabularx}
\caption{control\_key 定义}
\end{table}

\subsection{状态发布}

\begin{table}[H]
\centering
\begin{tabularx}{\textwidth}{lXl}
\toprule
\textbf{话题} & \textbf{内容} & \textbf{频率} \\
\midrule
/zit6/state/pos & NED 位姿 $[x,y,z,roll,pitch,yaw]$ & 30Hz \\
/zit6/state/vel & FRD 速度 $[u,v,w,p,q,r]$ & 50Hz \\
/zit6/state/thr & 归一化力/力矩 $[Fx,Fy,Fz,Mroll,Mpitch,Myaw]$ & 30Hz \\
/zit6/state/status & 系统状态（ARM、层级、INS） & 10Hz \\
/zit6/state/zithbt & 节点心跳 & 1Hz \\
/zit6/log & ROS2 日志 & 事件驱动 \\
\bottomrule
\end{tabularx}
\caption{状态发布话题}
\end{table}

\subsection{订阅指令}

\begin{itemize}
    \item \texttt{/zit6/cmd/setpoint} — 6-DOF 控制设定点（ZitSetpoint）；Roll/Pitch 字段仅作兼容/观测载荷，始终旁路
    \item \texttt{/zit6/cmd/agxhbt} — 上位机心跳（UInt32，用于 ARM）
    \item \texttt{/zit6/cmd/ins} — 惯导指令（UInt8）
    \item \texttt{/zit6/cmd/servo} — 舵机角度（Float32）
    \item \texttt{/zit6/cmd/light} — LED 灯控（UInt8）
    \item \texttt{/zit6/sim/nav} — SITL 仿真导航（Float32MultiArray[12]）
\end{itemize}

\subsection{参数服务}

\begin{itemize}
    \item \texttt{/zit6/update\_params} — JSON 格式实时更新参数
    \item \texttt{/zit6/get\_params} — 按路径查询参数
\end{itemize}

% ============================================================================
\section{配置系统}

\subsection{config.json $\rightarrow$ C++ 结构体}

配置通过 \texttt{gen\_config.py} 脚本生成 C++ 头文件：

\begin{lstlisting}[caption={配置结构体示例}]
struct AxisConfig {
    float pos_kp, pos_ki, pos_kd;
    float pos_i_limit, pos_output_limit;
    float vel_kp, vel_ki, vel_kd;
    float vel_i_limit, vel_output_limit;
    float max_v, max_a, mass, drag;
};

struct ChassisConfig {
    bool planner_enabled;
    AxisConfig x, y, z, roll, pitch, yaw;
};
\end{lstlisting}

运行时通过 \texttt{/zit6/update\_params} 服务实时重载。

\subsection{全局配置实例}

\begin{lstlisting}[caption={全局配置访问}]
namespace auv::config {
    extern SystemConfig sys_config;  // 只读运行时配置
    extern const ParamMeta SYSTEM_PARAMS[];  // 参数注册表
    extern const size_t SYSTEM_PARAMS_COUNT;
}
\end{lstlisting}

% ============================================================================
\section{坐标系定义}

采用两套坐标系，\textbf{轴向定义不同}：

\begin{table}[H]
\centering
\begin{tabularx}{\textwidth}{llll}
\toprule
\textbf{坐标系} & \textbf{约定} & \textbf{X/Y/Z} & \textbf{原点} \\
\midrule
机体系 & FRD & 前/右/下 & AUV 重心 \\
世界系 & NED & 北/东/下 & INS 上电零点 \\
\bottomrule
\end{tabularx}
\caption{坐标系定义}
\end{table}

6-DOF 全量运动学变换：
\begin{equation}
\nu_{world} = R(\phi, \theta, \psi) \cdot \nu_{body}
\end{equation}

其中 $R(\eta)$ 为 $6\times6$ 分块矩阵：
\begin{equation}
R(\eta) = \begin{bmatrix}
R_{ZYX}(\phi,\theta,\psi) & 0_{3\times3} \\
0_{3\times3} & T(\phi,\theta)
\end{bmatrix}
\end{equation}

% ============================================================================
\section{仿真支持}

\subsection{SITL（Software-in-the-Loop）}
\begin{itemize}
    \item 通过 \texttt{/zit6/sim/nav} 话题接收仿真导航数据
    \item 队列长度 3，消费端 drain 到最新帧
    \item 可在无硬件条件下测试控制逻辑
\end{itemize}

\subsection{HITL（Hardware-in-the-Loop）}
\begin{itemize}
    \item \texttt{HitlSimulator} 类在固件内部模拟 AUV 动力学
    \item 每步积分推进器推力，更新位姿
    \item 用于推进器分配和 PID 参数调优
\end{itemize}

仿真模式切换通过 \texttt{config.json} 中的 \texttt{simulation} 字段控制。

\end{document}

